\chapter{Background}
\label{ch:background}

\section{Signature Correctness and EUF-CMA Security}

A signature scheme is a triple of probabilistic algorithms
\[
  (\mathsf{KeyGen}, \mathsf{Sign}, \mathsf{Verify}).
\]
At security parameter \(\lambda\), key generation samples
\((pk,sk)\leftarrow\mathsf{KeyGen}(1^\lambda)\), signing samples
\(\sigma\leftarrow\mathsf{Sign}(sk,m,ctx)\), and verification returns a bit
\(\mathsf{Verify}(pk,m,ctx,\sigma)\in\{0,1\}\).  In the HAETAE development the
security parameter is represented indirectly by a fixed parameter mode and by
the symbolic hardness and loss terms attached to that mode, but the mathematical
reading is the usual asymptotic one.

Correctness is the assertion that honestly signed messages verify except with
negligible probability:
\[
  \Pr\!\left[
    \mathsf{Verify}(pk,m,ctx,\sigma)=1
    \;\middle|\;
    (pk,sk)\leftarrow\mathsf{KeyGen},\;
    \sigma\leftarrow\mathsf{Sign}(sk,m,ctx)
  \right]\ge 1-\varepsilon_{\mathsf{corr}}.
\]
Perfect correctness is the special case \(\varepsilon_{\mathsf{corr}}=0\).
For \eufcma{} security, an adversary receives \(pk\), random-oracle access, and
adaptive signing-oracle access.  If \(Q_S\) is the set of signed pairs
\((m,ctx)\), its success event is
\[
  \mathsf{Forge}(A) \equiv
  \mathsf{Verify}(pk,m^*,ctx^*,\sigma^*)=1
  \land (m^*,ctx^*)\notin Q_S .
\]
The advantage is
\[
  \Adv_{\Pi}^{\eufcma}(A)=\Pr[\mathsf{Forge}(A)].
\]
The proof objective is not to show that this probability is zero, but to bound
it by hardness terms and explicitly counted statistical losses.  A reductionist
statement has the shape
\[
  \Adv_{\Pi}^{\eufcma}(A)
  \le
  \sum_i \Adv_{\mathcal{P}_i}(B_i)
  + \sum_j L_j(q_H,q_S,\lambda),
\]
where, in a conventional paper proof, the \(B_i\) are adversaries constructed
against underlying mathematical problems and the \(L_j\) are loss terms arising
from collisions, prequeries, sampling changes, and rejection events.  In the
final EasyCrypt theorem checked in this dissertation, the corresponding MLWE and
Module-SIS solvers are declared reduction modules supplied through explicit
antecedents rather than constructed inside the top-level composition lemma.

The checked development models these notions in \code{Sig\_ROM.ec}.  This file
contains the generic interfaces for random oracles, schemes, signing oracles,
chosen-message adversaries, and no-message adversaries.  Keeping these games
generic is important: it prevents \haetae{}-specific proof assumptions from
being smuggled into the security definition.

\section{The Random Oracle Model}

The proof is a random-oracle-model proof.  A random oracle is an ideal sampled
function
\[
  H : \mathcal{X}\to\mathcal{Y}
\]
such that the values \(H(x)\) are independent samples from the prescribed output
distribution, subject only to consistency on repeated inputs.  In a lazy
implementation, the oracle maintains a finite partial map
\[
  T : \mathcal{X} \rightharpoonup \mathcal{Y}.
\]
On query \(x\), the oracle returns \(T(x)\) if \(x\in\operatorname{dom}(T)\);
otherwise it samples \(y\leftarrow D_x\), stores \(T[x\mapsto y]\), and returns
\(y\).  This lazy semantics is distributionally equivalent to sampling the whole
function in advance, but it exposes the finite set of adversarially observed
queries.

ROM proofs often program the oracle: a simulator chooses a point \(x_0\) and
sets \(H(x_0)\) to a value needed by the reduction.  Programming is sound only
when \(x_0\) is fresh:
\[
  x_0\notin Q_H,
\]
where \(Q_H\) is the adversary's prior query set.  If \(x_0\) has min-entropy at
least \(\log_2 N\) against the adversary's view and \(|Q_H|\le q_H\), then
\[
  \Pr[x_0\in Q_H]\le q_H/N.
\]
This inequality is the mathematical source of many counted ROM losses.  The
machine-checked development does not allow this argument to remain implicit:
the query log, budget counter, freshness predicate, and bad event are explicit
program state.

In the machine-checked proof, the random oracle is lazy and stateful.  Its map,
query log, bad flags, and counters are explicit program variables.  This makes
freshness and query-budget arguments checkable.

\section{Fiat-Shamir with Aborts}

\haetae{} follows the general pattern of Fiat-Shamir signatures with rejection
sampling.  The abstract pattern starts from an identification protocol with a
commitment, a challenge, and a response.  The Fiat-Shamir transform replaces the
verifier's challenge by a hash value
\[
  c = H(\mathsf{commitment},m,ctx).
\]
Security then depends on two interacting principles.

First, special soundness says that two accepting transcripts sharing the same
commitment but carrying different challenges reveal a witness or a solution to a
hard algebraic problem.  In lattice signatures this solution is typically a
short relation, hence the appearance of Module-SIS-style assumptions.

Second, the signing algorithm must produce transcripts with the same
adversary-visible distribution as the honest protocol would produce.  Rejection
sampling enforces this distributional property by discarding responses whose
shape would leak information about the secret key.  Mathematically, if
\(D_{\mathsf{real}}\) is the real signing distribution and
\(D_{\mathsf{sim}}\) is the simulated transcript distribution, the proof needs
a bound on statistical distance:
\[
  \Delta(D_{\mathsf{real}},D_{\mathsf{sim}})
  =
  \frac12\sum_x
  \left|\Pr[D_{\mathsf{real}}=x]-\Pr[D_{\mathsf{sim}}=x]\right|.
\]
The EasyCrypt development decomposes this distance into named sampling,
rejection, and clean-conditioned lifting lemmas rather than treating it as one
informal statement.

The word ``aborts'' refers to the rejection step.  It is not merely an
implementation branch: it is a proof device that shapes the output
distribution.  A correct provable-security proof must track both the probability
of aborting and the distribution conditioned on not aborting.

\section{Lattice Assumptions and Reduction Targets}

HAETAE's formal security proof ultimately reduces forgeries to structured
lattice problems.  At the level needed for the dissertation, the key point is
not the exact estimator cost of those problems, but the way the reduction
connects transcript algebra to assumption games.

Module-LWE is a pseudorandomness assumption.  In a typical module form, samples
have the shape
\[
  (A,\; A s + e)
\]
over a module \(R_q^k\), where \(s\) and \(e\) are short or noisy.  The
assumption says that these samples are computationally indistinguishable from
uniform samples of the same type.

Module-SIS is a short-relation assumption.  Given a public matrix \(A\), the
goal is to find a nonzero short vector \(z\) such that
\[
  A z = 0 \pmod q
\]
or, in inhomogeneous variants, \(A z = t \pmod q\).  Fiat-Shamir extraction
typically produces such a relation from two accepting transcripts with the same
commitment and different challenges.

The checked theorem keeps the hardness terms symbolic.  This is the right proof
engineering boundary: conditional on the declared hop and reduction premises,
EasyCrypt checks that the formal game bound composes the declared reduction
modules, symbolic assumptions, and declared probability losses.  It does not
perform independent lattice-estimator analysis for HAETAE parameter levels.

\section{EasyCrypt}

\easycrypt{} supports game-based cryptographic proofs using probabilistic
programs, module systems, Hoare logic, probabilistic Hoare logic, and relational
equivalence.  The \haetae{} development uses these features to express game
hops, oracle wrappers, state-preservation invariants, and concrete probability
bounds.

The mathematical interpretation of the main judgment forms is as follows.
A Hoare judgment proves that a program transforms states satisfying a
precondition into states satisfying a postcondition.  A probabilistic Hoare
judgment proves an upper or lower bound on the probability of a postcondition.
An equivalence judgment proves a coupling between two probabilistic programs,
usually the formal counterpart of a game hop.  Losslessness judgments prove that
procedures terminate with probability one, which is required before probabilities
can be compared without hidden nontermination mass.

For a provable-security proof, these judgments correspond to familiar paper
proof obligations:
\[
\begin{array}{rcl}
\text{state invariant} &\leftrightarrow& \text{Hoare judgment},\\
\text{bad-event probability} &\leftrightarrow& \text{probabilistic Hoare judgment},\\
\text{game indistinguishability} &\leftrightarrow& \text{equivalence judgment},\\
\text{well-defined experiment} &\leftrightarrow& \text{losslessness}.
\end{array}
\]
The value of the machine-checked proof is that these obligations are not
collapsed into prose.  Each must be represented as a typed statement over the
actual programs used by the security theorem.
